% GENERATED FILE. Edit canonical lessons, then run npm run generate:books.
% canonical-source-hash: fnv1a64:16afb4c6ad106259
% canonical-lessons: FR-C420-turbulence, FR-C420-escale, FR-C420-decaler, FR-C420-decalage, FR-C420-retarder

\chapter{Escale et retard}
\label{ch:fr-escale-et-retard}
\hlchaptermodality{420}

\begin{chapteropening}
\textbf{By the end of this chapter:} \emph{I can talk about bumpy flights and changed times: turbulence, the stopover, jet lag, moving a time and delaying.}

Complete the last lesson of chapter 420: I can talk about bumpy flights and changed times: turbulence, the stopover, jet lag, moving a time and delaying.
\end{chapteropening}

\section[la turbulence]{la turbulence — the turbulence}

\label{lesson:FR-C420-turbulence}



\begin{quote}
Before the new one: say the French for to land, then the French for to fly over.
\end{quote}

\subsection*{You'll want to know: la turbulence}
\textbf{la turbulence} — \textquotedblleft{}the turbulence\textquotedblright{}.

\begin{quote}
Pendant les turbulences, tout le monde reste assis.
\end{quote}

\begin{quote}
\emph{During the turbulence everyone stays seated.}
\end{quote}

Often plural: \emph{des turbulences}.

\begin{cousinweb}
\emph{turbulence} is a learned borrowing from Late Latin \emph{turbulentia}, from \emph{turba} 'disorder, crowd' (Lewis \& Short). EtymDB does not list it.
\end{cousinweb}

\begin{grammarlens}[title={What you've built}]
One more word for telling, explaining and describing at B1.
\end{grammarlens}

\subsection*{Your turn}
\begin{itemize}
\raggedright
  \item \emph{Say it:} \emph{la turbulence}
  \item \emph{Say it:} \emph{la turbulence}, once more
  \item \emph{Say it:} a sentence of your own with \emph{la turbulence}
  \item \emph{Recall:} say the French for to land, then the French for to fly over, then say \emph{la turbulence} again
\end{itemize}

\begin{tcolorbox}[breakable,colback=teal!4,colframe=teal!35!black,title={Before you move on}]
What does la turbulence mean? (\textquotedblleft{}The turbulence\textquotedblright{}.) Say it once more.
\end{tcolorbox}
\section[l'escale]{l'escale — the stopover, the port of call}

\label{lesson:FR-C420-escale}



\begin{quote}
Before the new one: say the French for to fly over, then the French for the turbulence.
\end{quote}

\subsection*{You'll want to know: l'escale}
\textbf{l'escale} — \textquotedblleft{}the stopover, the port of call\textquotedblright{}.

\begin{quote}
Nous avons fait une escale de trois heures.
\end{quote}

\begin{quote}
\emph{We had a three-hour stopover.}
\end{quote}

A feminine noun. \emph{Un vol sans escale} — a non-stop flight.

\begin{cousinweb}
\emph{escale} was borrowed from Italian \emph{scala} 'landing stage', from Latin \emph{scāla} 'ladder, steps' (EtymDB): the steps where a ship let people off.
\end{cousinweb}

\begin{grammarlens}[title={What you've built}]
One more word for telling, explaining and describing at B1.
\end{grammarlens}

\subsection*{Your turn}
\begin{itemize}
\raggedright
  \item \emph{Say it:} \emph{l'escale}
  \item \emph{Say it:} \emph{l'escale}, once more
  \item \emph{Say it:} a sentence of your own with \emph{l'escale}
  \item \emph{Recall:} say the French for to fly over, then the French for the turbulence, then say \emph{l'escale} again
\end{itemize}

\begin{tcolorbox}[breakable,colback=teal!4,colframe=teal!35!black,title={Before you move on}]
What does l'escale mean? (\textquotedblleft{}The stopover, the port of call\textquotedblright{}.) Say it once more.
\end{tcolorbox}
\section[décaler]{décaler — to shift, to move (a time)}

\label{lesson:FR-C420-decaler}



\begin{quote}
Before the new one: say the French for the turbulence, then the French for the stopover.
\end{quote}

\subsection*{You'll want to know: décaler}
\textbf{décaler} — \textquotedblleft{}to shift, to move (a time)\textquotedblright{}.

\begin{quote}
La compagnie a décalé le départ d'une heure.
\end{quote}

\begin{quote}
\emph{The airline moved the departure by an hour.}
\end{quote}

\emph{Décaler un rendez-vous} — to move an appointment.

\begin{cousinweb}
\emph{décaler} is built from \emph{dé-} and \emph{caler} 'to wedge': to take out the wedge and move something along. The origin of \emph{cale} 'wedge' is usually traced to German \emph{Keil}.
\end{cousinweb}

\begin{grammarlens}[title={What you've built}]
One more word for telling, explaining and describing at B1.
\end{grammarlens}

\subsection*{Your turn}
\begin{itemize}
\raggedright
  \item \emph{Say it:} \emph{décaler}
  \item \emph{Say it:} \emph{décaler}, once more
  \item \emph{Say it:} a sentence of your own with \emph{décaler}
  \item \emph{Recall:} say the French for the turbulence, then the French for the stopover, then say \emph{décaler} again
\end{itemize}

\begin{tcolorbox}[breakable,colback=teal!4,colframe=teal!35!black,title={Before you move on}]
What does décaler mean? (\textquotedblleft{}To shift, to move (a time)\textquotedblright{}.) Say it once more.
\end{tcolorbox}
\section[le décalage]{le décalage — the gap, the shift; jet lag}

\label{lesson:FR-C420-decalage}



\begin{quote}
Before the new one: say the French for the stopover, then the French for to shift.
\end{quote}

\subsection*{You'll want to know: le décalage}
\textbf{le décalage} — \textquotedblleft{}the gap, the shift; jet lag\textquotedblright{}.

\begin{quote}
Avec le décalage horaire, il dort mal pendant trois jours.
\end{quote}

\begin{quote}
\emph{With jet lag he sleeps badly for three days.}
\end{quote}

\emph{Le décalage horaire} — jet lag, the time difference.

\begin{cousinweb}
\emph{décalage} is built on \emph{décaler} 'to shift' with \emph{-age} (EtymDB).
\end{cousinweb}

\begin{grammarlens}[title={What you've built}]
One more word for telling, explaining and describing at B1.
\end{grammarlens}

\subsection*{Your turn}
\begin{itemize}
\raggedright
  \item \emph{Say it:} \emph{le décalage}
  \item \emph{Say it:} \emph{le décalage}, once more
  \item \emph{Say it:} a sentence of your own with \emph{le décalage}
  \item \emph{Recall:} say the French for the stopover, then the French for to shift, then say \emph{le décalage} again
\end{itemize}

\begin{tcolorbox}[breakable,colback=teal!4,colframe=teal!35!black,title={Before you move on}]
What does le décalage mean? (\textquotedblleft{}The gap, the shift; jet lag\textquotedblright{}.) Say it once more.
\end{tcolorbox}
\section[retarder]{retarder — to delay, to hold up}

\label{lesson:FR-C420-retarder}



\begin{quote}
Before the new one: say the French for to shift, then the French for the gap.
\end{quote}

\subsection*{You'll want to know: retarder}
\textbf{retarder} — \textquotedblleft{}to delay, to hold up\textquotedblright{}.

\begin{quote}
La neige a retardé tous les vols.
\end{quote}

\begin{quote}
\emph{The snow delayed all the flights.}
\end{quote}

\emph{Ma montre retarde} — my watch is slow.

\begin{cousinweb}
\emph{retarder} is a learned borrowing from Latin \emph{retardāre} 'to hold back', from \emph{tardus} 'slow' (Lewis \& Short).
\end{cousinweb}

\begin{grammarlens}[title={What you've built}]
One more word for telling, explaining and describing at B1.
\end{grammarlens}

\subsection*{Your turn}
\begin{itemize}
\raggedright
  \item \emph{Say it:} \emph{retarder}
  \item \emph{Say it:} \emph{retarder}, once more
  \item \emph{Say it:} a sentence of your own with \emph{retarder}
  \item \emph{Recall:} say the French for to shift, then the French for the gap, then say \emph{retarder} again
\end{itemize}

\begin{tcolorbox}[breakable,colback=teal!4,colframe=teal!35!black,title={Before you move on}]
What does retarder mean? (\textquotedblleft{}To delay, to hold up\textquotedblright{}.) Say it once more.
\end{tcolorbox}
